\documentclass[10pt,a4paper]{amsart}
\usepackage[margin=19mm]{geometry}
\usepackage{amsmath,amssymb,mathtools}
\usepackage[scr=rsfs,cal=euler]{mathalfa}
\usepackage{xfrac}
\usepackage[unicode,hidelinks,bookmarks=false]{hyperref}
\newcommand{\ie}{i.e.\ }
\newcommand{\cf}{cf.\ }
\newcommand{\resp}{respectively\ }
\newcommand{\Subsection}{\subsection}
\providecommand{\nobreakdash}{\nobreak\hbox{-}}
\setlength{\emergencystretch}{2em}
\multlinegap0pt
\usepackage{bm}
\usepackage{bbm}
\usepackage{dsfont}
\usepackage{xspace}
\usepackage{cancel}
\usepackage{mathtools}
\usepackage{amsthm}
\makeatletter
\@ifundefined{proposition}{\newtheorem{proposition}{Proposition}}{}
\@ifundefined{theorem}{\newtheorem{theorem}{Theorem}}{}
\makeatother
\title[Principal frequency of clamped plates on RCD(0,N) spaces: sh]{Principal frequency of clamped plates on RCD(0,N) spaces: sharpness, rigidity and stability}
\author{Alexandru Krist\'aly, Andrea Mondino}
\date{Source: arXiv:2409.04337v1 (CC BY 4.0)}
\begin{document}
\maketitle

\noindent\emph{Source and licence.} Principal frequency of clamped plates on RCD(0,N) spaces: sharpness, rigidity and stability. Version arXiv:2409.04337v1 (CC BY 4.0), distributed under the Creative Commons Attribution 4.0 International licence, as recorded in the arXiv licence metadata for this exact version and shown on the abstract page https://arxiv.org/abs/2409.04337v1. Full rights notice: licenses/arxiv-2409.04337-CC-BY-4.0.txt.\par\smallskip

\noindent\textit{Editorial scope.} Counted unit: the Theorem whose source label is \texttt{talenti-result-thm} in the article (source0.tex lines 641--650, source label \texttt{talenti-result-thm}), delivered together with the article's own complete original proof of it (source0.tex lines 651--677, one \texttt{proof} environment of 27 lines). No mathematics is written, repaired or re-derived: statement and proof are the article's own text line for line. Required context retained uncounted: equimeasure (lines 402--404); definition-F-pm (lines 539--541); simult=0 (lines 543--545); v-definit (lines 589--591); w-definit (lines 593--595); talenti-result-0 (lines 598--606); def-1-laplace (lines 634--636). Every definition, display and label of those lines is retained unchanged. Omitted from this package and not redistributed: 1--401, 405--538, 542--542, 546--588, 592--592, 596--597, 607--633, 637--640, 678--1362. Nothing inside a retained excerpt is omitted or altered: every retained body is delivered exactly as the article's lines are written, so no omission receipt of any kind accompanies this package. Citations: the retained excerpts carry no citation command at all, so no bibliography entry of the article is transcribed and no reference is redistributed. Reference adaptation: none was needed and none was made; every reference inside the delivered unit resolves inside it, so no unresolved reference or citation is produced. Printed slips kept literal: no printed word, symbol, hypothesis or number of the counted statement or of its proof was altered. Layout and class: the article's own class is \texttt{amsart} with packages including mathrsfs, srcltx, longtable, amsmath, xcolor, hyperref, geometry, bbm, dsfont; the wrapper is the lane's pinned \texttt{amsart} editorial wrapper at 10pt on A4, which restates the theorem-like environments the retained text needs and adds the article's own macro definitions that the retained text needs (none), with extra wrapper packages (bm, bbm, dsfont, xspace, cancel, mathtools). The delivered numbering is the unit's own. Assets: no retained span contains a figure environment, a graphics-inclusion command, a PDF-page inclusion, a table or a TikZ picture; no image file is copied into this package, the package declares no asset dependency and no image file is redistributed with it. Licence: version arXiv:2409.04337v1 of this article is distributed under the Creative Commons Attribution 4.0 International licence, as recorded in the arXiv licence metadata for this exact version and shown on the abstract page https://arxiv.org/abs/2409.04337v1. A full rights notice is delivered at licenses/arxiv-2409.04337-CC-BY-4.0.txt. Characters: every retained excerpt is pure ASCII, so the delivered engine, XeLaTeX with its default Latin Modern text font, reports no missing character.
		\begin{equation}\label{equimeasure}
		\|u\|_{L^p(\Omega,{\sf m})}=\|u^*\|_{L^p(\Omega^*,\sigma_N)},  \quad \text{for every } p\geq 1,
	\end{equation} 

\begin{equation}\label{definition-F-pm}
	F_+(s)=(\Delta u)^\#_-(s)-(\Delta u)^\#_+({\sf m}(\tilde \Omega)-s)\ \ {\rm and}\ \ F_-(s)=-F_+({\sf m}(\tilde \Omega)-s),\ \ s\in [0,{\sf m}(\tilde \Omega)].
\end{equation}

\begin{equation}\label{simult=0}
	(\Delta u)^\#_-(s)(\Delta u)^\#_+({\sf m}(\tilde \Omega)-s)=0,\ \ s\in [0,{\sf m}(\tilde \Omega)].
\end{equation}

\begin{equation}\label{v-definit}
	V_+(x)=\frac{1}{N\omega_N}\int_{x}^a \rho^{1-N}\left(\int_0^{\omega_N \rho^N}F_+(s){\rm d}s\right){\rm d}\rho,\ \ x\in [0,a],
	\end{equation}

	\begin{equation}\label{w-definit}
		V_-(x)=\frac{1}{N\omega_N}\int_{x}^b \rho^{1-N}\left(\int_0^{\omega_N \rho^N}F_-(s){\rm d}s\right){\rm d}\rho,\ \ x\in [0,b],
				\end{equation}

		\begin{proposition} \label{talenti-result-0}
			Let  $V_+$ and $V_-$ from {\rm (\ref{v-definit})} and {\rm (\ref{w-definit})}, respectively. Then
			\begin{equation}\label{1-compar-0}
				{\sf AVR}_ {\sf m}^\frac{2}{N}	u_+^*\leq V_+\ \ {in}\ \ [0,a];
			\end{equation}
			\begin{equation}\label{2-compar-0}
				{\sf AVR}_ {\sf m}^\frac{2}{N}	u_-^*\leq V_-\ \ {in}\ \ [0,b].
			\end{equation}
		\end{proposition}

		\begin{equation}\label{def-1-laplace}
			\Delta_{0,N} V(r)=V''(r)+\frac{N-1}{r}V'(r),\ \ r>0.
		\end{equation}

		\begin{theorem} \label{talenti-result-thm}
			Let  $V_+$ and $V_-$ from {\rm (\ref{v-definit})} and {\rm (\ref{w-definit})}, respectively. Then
			\begin{equation}\label{u-v-w}
				{\sf AVR}_ {\sf m}^\frac{4}{N}	\int_{\tilde \Omega}u^2 \, {\rm d}{\sf m}\leq \int_{0}^aV_+^2 \, {\rm d}\sigma_N+\int_{0}^b V_-^2 \, {\rm d}\sigma_N,
			\end{equation}
			and
			\begin{equation}\label{laplace-hasonlitas}
				\int_{\tilde \Omega} (\Delta u)^2 \, {\rm d}{\sf m}=\int_{0}^a (\Delta_{0,N}V_+)^2 \, {\rm d}\sigma_N+ \int_{0}^b (\Delta_{0,N}V_-)^2 \, {\rm d}\sigma_N.
			\end{equation}
		\end{theorem}

		\begin{proof}
			The proof of 	\eqref{u-v-w} follows by Proposition \ref{talenti-result-0} combined with \eqref{equimeasure}; indeed, one has that 
			\begin{eqnarray*}
				{\sf AVR}_ {\sf m}^\frac{4}{N}	\int_{\tilde \Omega}u^2 {\rm d}{\sf m}&=&{\sf AVR}_ {\sf m}^\frac{4}{N}	\int_{\Omega_+}u_+^2 {\rm d}{\sf m}+{\sf AVR}_ {\sf m}^\frac{4}{N}	\int_{ \Omega_-}u_-^2 {\rm d}{\sf m}\\&=&{\sf AVR}_ {\sf m}^\frac{4}{N}	\int_{0}^a (u_+^*)^2 {\rm d}\sigma_N+{\sf AVR}_ {\sf m}^\frac{4}{N}	\int_{ 0}^b(u_-^*)^2 {\rm d}\sigma_N\\&\leq&\int_{0}^aV_+^2 {\rm d}\sigma_N+\int_{0}^b V_-^2 {\rm d}\sigma_N.
			\end{eqnarray*}
			Let us focus now on \eqref{laplace-hasonlitas}. First, due to \eqref{simult=0} and \eqref{equimeasure}, we observe that
			\begin{align*}
				\int_{0}^{{\sf m}(\tilde \Omega)} F^2_+(s){\rm d}s
				&=\int_{0}^{{\sf m}(\tilde \Omega)} \left[(\Delta u)^\#_-(s)^2+(\Delta u)^\#_+({\sf m}(\tilde \Omega)-s)^2-2(\Delta u)^\#_-(s)(\Delta u)^\#_+({\sf m}(\tilde \Omega)-s)\right]{\rm d}s
				\\
				&=\int_{0}^{\omega_N \tilde R^N} \left[(\Delta u)^\#_-(s)^2+(\Delta u)^\#_+(s)^2\right]{\rm d}s\\&=\int_{0}^{\tilde R} \left[(\Delta u)^*_-(t)^2+(\Delta u)^*_+(t)^2\right]{\rm d}\sigma_N(t)	=\int_{\tilde \Omega} \left[(\Delta u)_-^2+(\Delta u)^2_+\right]{\rm d}{\sf m}
				\\&=\int_{\tilde \Omega} (\Delta u)^2{\rm d}{\sf m}.
			\end{align*}
			On the other hand, by using \eqref{def-1-laplace}, 
			a simple computation shows that
			\begin{equation}\label{laplace-1}
				\Delta_{0,N} V_+(r)=-F_+(\omega_N r^N),\ \ r\in (0,a],
			\end{equation}
			\begin{equation}\label{laplace-2}
				\Delta_{0,N} V_-(r)=-F_-(\omega_N r^N),\ \ r\in (0,b].
			\end{equation}
			Therefore, by \eqref{laplace-1},  \eqref{laplace-2} and  \eqref{definition-F-pm}, combined with the fact that $a^N+b^N=\tilde R^N$, one has that
			\begin{eqnarray*}
				\int_{0}^a (\Delta_{0,N}V_+)^2{\rm d}\sigma_N+ \int_{0}^b (\Delta_{0,N}V_-)^2{\rm d}\sigma_N&=&\int_{0}^a F_+^2(\omega_N r^N){\rm d}\sigma_N(r)+ \int_{0}^b F^2_-(\omega_N r^N){\rm d}\sigma_N(r)\\&=& \int_{0}^{\omega_N a^N} F_+^2(s){\rm d}s+ \int_{0}^{\omega_N b^N} F_-^2(s){\rm d}s\\&=&\int_{0}^{\omega_N \tilde R^N} F_+^2(s){\rm d}s.
			\end{eqnarray*}
			Since ${\sf m}(\tilde \Omega)=\omega_N \tilde R^N$, by combining the above findings, relation  \eqref{laplace-hasonlitas} yields at once. 
		\end{proof}

\end{document}
